\chapter[Motivating Examples]{Motivating Examples}\label{ch:sec:motivating-examples}

This book develops a mathematical framework---weighted incidence
structures, their forced Hilbert-space geometry, and the Universal
Optimality Defect---that unifies the analysis of encryption, privacy,
security, and information measures.  Before building the theory, we
present four concrete problems that it solves.

\section{AES and perfect secrecy}\label{ch:sec:aes-and-perfect-secrecy}

\paragraph{Main running example: the AES S-box.}
Throughout the book we return to one concrete encryption system to
illustrate every new concept as it is introduced.  The Advanced
Encryption Standard (AES) encrypts a 128-bit plaintext $M$ using a
128-bit key $K$ to produce a 128-bit ciphertext $C$.
Every pair $(M,K)$ maps to a distinct ciphertext.  After observing
$C$, an attacker knows that every plaintext $M$ is compatible with
exactly one key---there is no posterior ambiguity.  The cipher is
perfectly secret in the sense of Shannon.

In our framework this is the simplest case: the WIS has $n_{mc}=1$ for
every triple $(M,C)$.  The Universal Optimality Defect is $0$, and the
classical theory of this book recovers perfect secrecy as the trivial
case of a general algebraic structure.

\paragraph{Additional running examples.}
Three further problems complement the main example and illustrate the
breadth of the framework.

\section{Browser fingerprinting and anonymity}\label{ch:sec:browser-fingerprinting-and-anonymity}

\textbf{Browser fingerprinting and anonymity.}
When a user visits a website, their browser reveals dozens of
attributes---screen resolution, installed fonts, time zone, user
agent, installed plugins, HTTP headers.  The
Panopticlick~\cite{eckersley2010unique} study found that $83.6\%$ of
desktop browsers had a unique fingerprint among those tested.
This is a WIS with a large message space (browsers), a large
ciphertext space (fingerprints), and a highly irregular incidence
structure (most browsers produce unique fingerprints, but some are
shared by many).  The Universal Optimality Defect measures how far
this system is from the uniform posterior ideal: a high UOD means
the fingerprint is strongly identifying.

\section{Geo-indistinguishability and location privacy}\label{ch:sec:geo-indistinguishability-and-location-privacy}

\textbf{Geo-indistinguishability and location privacy.}
A user wishes to report their location to a service without revealing
it exactly.  The geo-indistinguishability model
\cite{andres2013geo} adds planar Laplacian noise to the true
coordinates: $C=M+\operatorname{Lap}(0,\varepsilon^{-1})^2$.  This
ensures that any two locations within distance $r$ produce
$\varepsilon r$-close ciphertext distributions.
The WIS is continuous rather than discrete, but the same principles
apply: the Laplacian mechanism is a universally optimal encoder in the
sense that the posterior is uniform on a closed disk centred at the
reported location.  The trade-off between privacy (large noise) and
utility (small noise) is captured by the UOD, which in this case
equals $2/\varepsilon^2$.

\section{Differential privacy and database queries}\label{ch:sec:differential-privacy-and-database-queries}

\textbf{Differential privacy and database queries.}
A statistical database answers queries about its rows.  Differential
privacy \cite{dwork2006differential} requires that adding or removing
a single row changes the output distribution by at most $\varepsilon$
multiplicatively.  This is equivalent to a WIS condition: every pair
of neighbouring databases must be incognisable.
The UOD measures the distinguishability of neighbouring databases.
The classical Laplace mechanism achieves $\mathcal D=2/\varepsilon^2$.
Our geometric bounds show that this distinguishability is the minimum
possible for any $\varepsilon$-differentially private mechanism,
providing a new proof of the optimality of the Laplace mechanism.

\paragraph{How to read this book.}
The AES S-box is our main running example and reappears in every
chapter to ground abstract concepts in concrete numbers.
The additional examples are used where they illuminate a point
particularly well---fingerprinting for piecewise-smooth functionals,
geo-indistinguishability for continuous limits, differential privacy
for optimality bounds.  When a new concept is introduced---WIS,
gauge symmetry, the Universal Optimality Defect, the Lipschitz or
quadratic bounds, the Bregman framework, the Three-Level Theorem---we
show how it applies first to the AES S-box and then to the additional
examples.
Chapter~\ref{ch:sec:applications-of-the-wis-framework} collects eight
extended applications with full mathematical detail.
